\documentclass[12pt]{article}
\usepackage[utf8]{inputenc}
\usepackage[T1]{fontenc}
\usepackage[margin=1in]{geometry}
\usepackage{setspace}
\usepackage{xcolor}
\usepackage[colorlinks=true, urlcolor=blue!60!black, linkcolor=blue!60!black]{hyperref}
\setlength{\parindent}{2em}
\setlength{\parskip}{1.25em}
\renewcommand{\baselinestretch}{1.0}
\begin{document}
\begin{center}
{\LARGE \textit{Louisiana v. Callais}, 608 U.S. \_\_\_ (2026)\par}
\vspace{0.8em}
{\large Case Brief\par}
\end{center}
\vspace{1em}

\section{Memory Jogger}

In \href{https://www.courtlistener.com/opinion/10852760/louisiana-v-callais-revisions-50426/}{\textit{Louisiana v. Callais}}, the Supreme Court held 6--3 that Louisiana's congressional map creating a second majority-Black district was an unconstitutional racial gerrymander under the Equal Protection Clause because Section~2 of the Voting Rights Act---reinterpreted to require a strong inference of intentional racial discrimination and to disentangle race from politics---did not require the second district and therefore supplied no compelling governmental interest.

\section{Detailed Case Facts}

Following the 2020 decennial census, Louisiana retained its allocation of six seats in the U.S. House of Representatives, with Black residents comprising roughly one-third of the state's population (\href{https://supreme.justia.com/cases/federal/us/608/24-109/}{608 U.S. at \_\_\_ (slip op., at 11)}; \href{https://www.scotusblog.com/2026/04/in-major-voting-rights-act-case-supreme-court-strikes-down-redistricting-map-challenged-as-racia/}{SCOTUSblog}). In 2022, the Louisiana Legislature enacted House Bill~1 (HB1), a congressional redistricting map that, like its 2013--2022 predecessor, contained only one majority-Black congressional district (District~2, anchored in New Orleans) and five majority-White districts (\href{https://supreme.justia.com/cases/federal/us/608/24-109/}{Justia}).

Black voters and civil rights plaintiffs (the \textit{Robinson} plaintiffs) sued in the U.S. District Court for the Middle District of Louisiana, alleging that HB1 diluted Black voting power in violation of Section~2 of the Voting Rights Act of 1965 (VRA), 52 U.S.C.~\S~10301, by ``packing'' Black voters into District~2 and ``cracking'' the remaining Black population across the other five districts (\href{https://supreme.justia.com/cases/federal/us/608/24-109/}{Justia}). After the Middle District concluded that HB1 likely violated Section~2 and the Fifth Circuit remanded for Louisiana to draw a remedial map or face a court-imposed plan, the Louisiana Legislature met in a special session in January 2024 and enacted Senate Bill~8 (SB8) (\href{https://www.scotusblog.com/2026/04/in-major-voting-rights-act-case-supreme-court-strikes-down-redistricting-map-challenged-as-racia/}{SCOTUSblog}; \href{https://supreme.justia.com/cases/federal/us/608/24-109/}{Justia}).

SB8 preserved majority-Black District~2 and added a second majority-Black district, District~6, configured to achieve a Black voting-age population exceeding 50\% (\href{https://supreme.justia.com/cases/federal/us/608/24-109/}{Justia}). Rather than adopting the \textit{Robinson} plaintiffs' illustrative maps---which connected urban Black communities in Baton Rouge and Lafayette with rural Black communities in northeast Louisiana and would have endangered Republican Representative Julia Letlow---the Legislature drew District~6 as a roughly 250-mile-long diagonal slash stretching from Shreveport in northwest Louisiana through Alexandria and Lafayette to Baton Rouge in the southeast (\href{https://supreme.justia.com/cases/federal/us/608/24-109/}{Justia}; \href{https://www.everycrsreport.com/reports/LSB11431.html}{CRS Report LSB11431}). Louisiana chose that configuration to satisfy the Middle District's VRA ruling while shielding key Republican incumbents: House Speaker Mike Johnson, House Majority Leader Steve Scalise, and Appropriations Committee member Julia Letlow (\href{https://supreme.justia.com/cases/federal/us/608/24-109/}{Justia}). Under SB8, former Representative Cleo Fields won election in District~6 in November 2024 (\href{https://www.scotusblog.com/2026/04/in-major-voting-rights-act-case-supreme-court-strikes-down-redistricting-map-challenged-as-racia/}{SCOTUSblog}).

\section{Procedural History}

In the initial VRA litigation, \textit{Robinson v. Ardoin}, 605 F.~Supp.~3d 759 (M.D. La. 2022), the Middle District of Louisiana preliminarily enjoined HB1 as likely violating Section~2. After the U.S. Supreme Court granted certiorari before judgment, held the case pending \textit{Allen v. Milligan}, 599 U.S. 1 (2023), and then dismissed the writ as improvidently granted, the Fifth Circuit held that the preliminary injunction ``was valid when it was issued'' but vacated it as no longer urgent and remanded to give the Legislature time to enact a new map, 86 F.4th 574 (5th Cir. 2023) (\href{https://supreme.justia.com/cases/federal/us/608/24-109/}{Justia}).

Shortly after Louisiana enacted SB8 in early 2024, twelve self-described ``non-African American'' voters led by Phillip Callais sued in the U.S. District Court for the Western District of Louisiana, claiming SB8 was an unconstitutional racial gerrymander under the Fourteenth Amendment's Equal Protection Clause (\href{https://www.scotusblog.com/2026/04/in-major-voting-rights-act-case-supreme-court-strikes-down-redistricting-map-challenged-as-racia/}{SCOTUSblog}; \href{https://supreme.justia.com/cases/federal/us/608/24-109/}{Justia}). The \textit{Robinson} plaintiffs intervened to defend SB8. A three-judge district court convened under 28 U.S.C.~\S~2284(a) held a three-day consolidated hearing and trial and ruled 2--1 (over Fifth Circuit Judge Carl Stewart's dissent) that SB8 violated the Equal Protection Clause, enjoining its use, \textit{Callais v. Landry}, 732 F.~Supp.~3d 574 (W.D. La. 2024) (\href{https://supreme.justia.com/cases/federal/us/608/24-109/}{Justia}).

The Supreme Court stayed the injunction in May 2024 for the 2024 election cycle and noted probable jurisdiction on direct appeal under 28 U.S.C.~\S~1253 on November 4, 2024, consolidating \textit{Louisiana v. Callais} (No.~24-109) with \textit{Robinson v. Callais} (No.~24-110) (\href{https://www.oyez.org/cases/2025/24-109}{Oyez}; \href{https://www.scotusblog.com/2026/04/in-major-voting-rights-act-case-supreme-court-strikes-down-redistricting-map-challenged-as-racia/}{SCOTUSblog}). Following initial oral argument on March 24, 2025, the Court on June 27, 2025---over Justice Thomas's dissent---restored the cases for reargument and ordered supplemental briefing on whether Louisiana's intentional creation of a second majority-minority district violated the Fourteenth or Fifteenth Amendments (\href{https://www.supremecourt.gov/opinions/24pdf/24-109_l53m.pdf}{606 U.S. 923 (2025)}; \href{https://supreme.justia.com/cases/federal/us/608/24-109/}{Justia}). By the October 15, 2025 reargument, both Louisiana and the United States (as amicus curiae) opposed SB8, leaving the \textit{Robinson} intervenors to defend the map (\href{https://www.scotusblog.com/2026/04/in-major-voting-rights-act-case-supreme-court-strikes-down-redistricting-map-challenged-as-racia/}{SCOTUSblog}). The Court issued its decision on April 29, 2026 (\href{https://www.supremecourt.gov/opinions/25pdf/24-109_21o3.pdf}{Slip Opinion}). On May 4, 2026, over Justice Jackson's dissent, the Court granted the \textit{Callais} appellees' application to issue the judgment forthwith ahead of Louisiana's scheduled May 16 primary (\href{https://www.supremecourt.gov/opinions/25pdf/25a1197_h31i.pdf}{No.~25A1197}), and on May 6, 2026, denied a motion to recall the judgment (\href{https://www.scotusblog.com/2026/05/court-gives-immediate-effect-to-voting-rights-act-decision/}{SCOTUSblog}).

\section{Judicial Votes}

\textbf{Vote:} 6--3 affirming the three-judge district court and remanding (\href{https://www.supremecourt.gov/opinions/25pdf/24-109_21o3.pdf}{Slip Opinion}; \href{https://www.scotusblog.com/cases/louisiana-v-callais-2/}{SCOTUSblog}).

\begin{itemize}
    \item \textbf{Majority (6):} Justice Samuel A. Alito, Jr. delivered the opinion of the Court, joined in full by Chief Justice John G. Roberts, Jr. and Justices Clarence Thomas, Neil M. Gorsuch, Brett M. Kavanaugh, and Amy Coney Barrett.
    \item \textbf{Concurrence (2):} Justice Clarence Thomas filed a concurring opinion, joined by Justice Neil M. Gorsuch.
    \item \textbf{Dissent (3):} Justice Elena Kagan filed a dissenting opinion, joined by Justices Sonia Sotomayor and Ketanji Brown Jackson.
\end{itemize}

\section{Holding}

\textbf{Judgment:} The judgment of the three-judge U.S. District Court for the Western District of Louisiana (732 F.~Supp.~3d 574) is affirmed, and the consolidated cases are remanded for further proceedings (\href{https://supreme.justia.com/cases/federal/us/608/24-109/}{Justia}).

\textbf{Rule Announced:} Compliance with Section~2 of the Voting Rights Act, \textit{as properly construed}, can constitute a compelling governmental interest satisfying Equal Protection strict scrutiny. However, to align with Section~2's text and Congress's enforcement power under Section~2 of the Fifteenth Amendment, Section~2 imposes vote-dilution liability ``only when the evidence supports a strong inference that the State intentionally drew its districts to afford minority voters less opportunity because of their race'' (\href{https://supreme.justia.com/cases/federal/us/608/24-109/}{Justia}). Updating the \textit{Thornburg v. Gingles}, 478 U.S. 30 (1986), framework, a Section~2 plaintiff must now:
\begin{enumerate}
    \item Under \textit{Gingles} Precondition~1, produce an illustrative map drawn \textit{without} using race as a districting criterion that satisfies all of the state's legitimate nonracial objectives, including traditional criteria and specified political or partisan goals (such as incumbency protection or partisan seat distribution);
    \item Under \textit{Gingles} Preconditions~2 and~3, demonstrate racially polarized voting through an analysis that controls for partisan affiliation; and
    \item Under the ``totality of circumstances,'' focus on current data showing an objective likelihood of present-day intentional racial discrimination in voting, rather than historical discrimination or ongoing effects of societal discrimination (\href{https://supreme.justia.com/cases/federal/us/608/24-109/}{Justia}; \href{https://www.everycrsreport.com/reports/LSB11431.html}{CRS Report LSB11431}).
\end{enumerate}

\textbf{Holding vs.\ Dicta:} The holding is that Louisiana's SB8 triggered strict scrutiny because its underlying goal was racial, that the \textit{Robinson} plaintiffs failed to establish a Section~2 violation under the updated \textit{Gingles} framework, and that SB8 is therefore an unconstitutional racial gerrymander unsupported by a compelling interest (\href{https://supreme.justia.com/cases/federal/us/608/24-109/}{Justia}). The Court's threshold determination that compliance with a properly construed Section~2 \textit{can} serve as a compelling interest, along with its computer-algorithm hypothetical illustrating when rejecting majority-minority maps without a legitimate explanation could raise a strong inference of intentional discrimination (\href{https://supreme.justia.com/cases/federal/us/608/24-109/}{Justia}), functions as explanatory framework and instructive dicta, since Louisiana's map failed because Section~2 did not require it in the first place. Likewise, the Court assumed without deciding that a legislative districting map constitutes a ``standard, practice, or procedure'' under Section~2(a) (\href{https://supreme.justia.com/cases/federal/us/608/24-109/}{Justia}).

\section{Analysis of Opinions}

\subsection{Majority Opinion (Justice Alito)}

Justice Alito framed the case as resolving a thirty-year tension left open from \textit{Miller v. Johnson} (1995) through \textit{Wisconsin Legislature v. Wisconsin Elections Commission} (2022): whether compliance with Section~2 of the VRA qualifies as a compelling interest justifying race-based redistricting under Equal Protection strict scrutiny (\href{https://supreme.justia.com/cases/federal/us/608/24-109/}{Justia}). Citing \textit{Students for Fair Admissions, Inc. v. President and Fellows of Harvard College}, 600 U.S. 181 (2023), Alito reiterated that the Constitution permits racial classifications only in the rarest circumstances---chiefly avoiding imminent prison safety risks or remediating specific, identified past discrimination. He held that compliance with Section~2 can be added to that short list, but only when Section~2 is properly interpreted to avoid colliding with the Constitution (\href{https://supreme.justia.com/cases/federal/us/608/24-109/}{Justia}).

Turning to the statutory text of 52 U.S.C.~\S~10301(b), Alito parsed ``less opportunity than other members of the electorate to ... elect representatives of their choice.'' Because Section~2 expressly disclaims any right to proportional representation, ``opportunity'' means a chance to elect, measured against a randomly selected comparator group of nonminority voters under the state's permissible districting criteria---including traditional criteria and partisan advantage (\href{https://supreme.justia.com/cases/federal/us/608/24-109/}{Justia}). Under \textit{City of Boerne v. Flores}, 521 U.S. 507 (1997), legislation enforcing the Fifteenth Amendment must be ``congruent and proportional'' to the Amendment's prohibition, which bars only purposeful racial discrimination (\textit{Reno v. Bossier Parish School Bd.}, 520 U.S. 471 (1997)). Accordingly, to stay within Congress's Fifteenth Amendment authority and invoke constitutional avoidance, Section~2 cannot punish mere disparate impact; it imposes liability only where the circumstances support a ``strong inference'' of intentional racial discrimination (\href{https://supreme.justia.com/cases/federal/us/608/24-109/}{Justia}; \href{https://www.everycrsreport.com/reports/LSB11431.html}{CRS Report LSB11431}).

Rather than formally overruling \textit{Gingles} or \textit{Allen v. Milligan}, Alito ``updated'' \textit{Gingles} based on four post-1986 developments: (1) dramatic social progress and equalized Black voter turnout, especially in the South (citing \textit{Shelby County v. Holder}, 570 U.S. 529 (2013)); (2) the transition from a one-party South to a competitive two-party system where race and party correlate strongly; (3) \textit{Rucho v. Common Cause}, 588 U.S. 684 (2019), which rendered partisan gerrymandering nonjusticiable and incentivized litigants to repackage partisan claims as Section~2 claims (citing \textit{Alexander v. South Carolina State Conference of the NAACP}, 602 U.S. 1 (2024)); and (4) modern computing power enabling plaintiffs to generate millions of alternative maps (\href{https://supreme.justia.com/cases/federal/us/608/24-109/}{Justia}). Distinguishing \textit{Allen v. Milligan} on the ground that Alabama neither raised Fourteenth Amendment racial-gerrymandering claims nor defended its map on partisan grounds, Alito applied the revised \textit{Gingles} test and held that the \textit{Robinson} plaintiffs failed at every step: their illustrative maps failed to protect Representative Letlow, their polarization studies did not control for party affiliation, and their totality-of-circumstances showing relied on historical discrimination rather than present-day intentional discrimination (\href{https://supreme.justia.com/cases/federal/us/608/24-109/}{Justia}).

\subsection{Concurring Opinion (Justice Thomas)}

Justice Thomas, joined by Justice Gorsuch, joined the Court's opinion in full because it ``should largely put an end to this `disastrous misadventure' in voting-rights jurisprudence'' that had turned Section~2 into a de facto mandate for racial proportionality repugnant to a color-blind Constitution (\href{https://supreme.justia.com/cases/federal/us/608/24-109/}{Justia}). Reiterating the textualist position he first articulated in \textit{Holder v. Hall}, 512 U.S. 874 (1994), and his dissent in \textit{Allen v. Milligan}, Thomas wrote that he would go further and hold that Section~2 does not apply to redistricting at all: its statutory terms---``voting qualification,'' ``prerequisite to voting,'' and ``standard, practice, or procedure''---reach only enactments regulating ballot access and vote counting, not a state's choice of district lines (\href{https://supreme.justia.com/cases/federal/us/608/24-109/}{Justia}; \href{https://www.everycrsreport.com/reports/LSB11431.html}{CRS Report LSB11431}).

\subsection{Dissenting Opinion (Justice Kagan)}

Justice Kagan, joined by Justices Sotomayor and Jackson, delivered a sweeping 48-page dissent (and read a summary from the bench, omitting ``respectfully'' from her closing ``I dissent'') (\href{https://www.scotusblog.com/2026/04/in-major-voting-rights-act-case-supreme-court-strikes-down-redistricting-map-challenged-as-racia/}{SCOTUSblog}). Opening with a stylized hypothetical of a state slicing a compact 90\% Black circular county into six pie pieces appended to 90\% White districts, Kagan demonstrated that the majority's new ``Callais requirements'' would immunize even textbook vote dilution whenever a state invokes a partisan gerrymandering goal (\href{https://supreme.justia.com/cases/federal/us/608/24-109/}{Justia}).

Kagan traced the 1982 amendments to Section~2, enacted by overwhelming bipartisan majorities specifically to abrogate \textit{City of Mobile v. Bolden}, 446 U.S. 55 (1980), and replace its intent requirement with the ``results'' test of \textit{White v. Regester}, 412 U.S. 755 (1973). By importing the Fourteenth Amendment racial-gerrymandering ``disentanglement'' burden from \textit{Cooper} and \textit{Alexander} into Section~2 vote-dilution suits, the majority converted Congress's effects test back into \textit{Bolden}'s virtually insurmountable intent test (\href{https://supreme.justia.com/cases/federal/us/608/24-109/}{Justia}). She further accused the majority of flouting ``superpowered'' statutory \textit{stare decisis} just three Terms after \textit{Allen v. Milligan} explicitly reaffirmed \textit{Gingles} and upheld Section~2's effects test under the Fifteenth Amendment. Viewing \textit{Callais} as the final step in a trilogy with \textit{Shelby County} (gutting Section~5) and \textit{Brnovich v. Democratic National Committee} (2021) (curtailing vote-denial claims), Kagan warned that the ruling leaves even longstanding majority-minority districts like Louisiana's District~2 vulnerable to elimination and threatens the largest reduction in minority representation since the end of Reconstruction (\href{https://supreme.justia.com/cases/federal/us/608/24-109/}{Justia}).

\section{Examples: Future Applications}

Applying \textit{Louisiana v. Callais} to future redistricting disputes produces the following outcomes:

\begin{enumerate}
    \item \textbf{Same-Side (Section~2 Claim Fails / Race-Based Remedy Invalidated): Cracking a Cohesive Black Urban District to Create a 6--0 Partisan Map.} A southern state legislature dismantles a longstanding, compact majority-Black congressional district in a major city and disperses its voters across four surrounding suburban districts, openly declaring its goal is to elect an all-Republican delegation. Because Black voters overwhelmingly support Democrats, no illustrative map can both preserve a majority-Black district and achieve the state's stated partisan seat distribution, nor can plaintiffs separate inter-party racial polarization from party preference. Under \textit{Callais}, the Section~2 claim fails at \textit{Gingles} Preconditions~1--3, and any court-ordered or voluntary race-based redraw would violate the Equal Protection Clause.

    \item \textbf{Same-Side (Section~2 Claim Fails): Failure to Protect Incumbents and Communities of Interest.} Plaintiffs challenge a state's congressional map (as in the post-\textit{Callais} stay order in \href{https://www.supremecourt.gov/opinions/25pdf/608us2r42_q8l1.pdf}{\textit{Allen v. Milligan}, 608 U.S. 511 (2026)}) by submitting an illustrative map with an additional majority-minority district that pairs two incumbents or splits a state-prioritized regional community of interest. Because the plaintiffs' map does not perform ``just as well'' on all of the state's constitutionally permissible districting criteria, the Section~2 claim fails under \textit{Callais}.

    \item \textbf{Opposite-Side (Section~2 Violation Found / Compelling Interest Established): Intra-Party Racial Exclusion in a One-Party Municipality.} In a heavily Democratic city where virtually all residents vote Democratic, White, Black, and Hispanic voters sharply divide along racial lines in Democratic primaries, and the city council redraws wards to split a compact Hispanic neighborhood despite computer-generated, race-blind maps satisfying every stated city criterion (compactness, incumbent protection, and precinct continuity) that keep the neighborhood intact. Here, polarization occurs \textit{within} a single party (controlling for partisanship), and the council's unexplained rejection of race-blind maps meeting all legitimate goals creates a strong inference of intentional racial discrimination under \textit{Callais}.

    \item \textbf{Opposite-Side (Section~2 Violation Found): Direct Evidence of Intentional Racial Cracking Accompanied by Race-Blind Alternatives.} During redistricting, leaked legislative communications and mapmaker notes show that officials deliberately targeted Black precincts because of race---not party---after a race-neutral algorithm produced multiple maps that equally achieved the legislature's partisan and incumbent-protection targets while preserving a majority-Black district. Because plaintiffs can satisfy all nonracial state goals with a race-blind alternative map and present current evidence of intentional racial discrimination, Section~2 liability attaches and supplies a compelling interest for a narrowly tailored remedy.

    \item \textbf{Fence-Sitter (Close Call): Direct Evidence of Racial Animus Without an Equally Effective Partisan Alternative Map.} Suppose discovery uncovers explicit emails from key legislators expressing hostility toward minority voters and a desire to eliminate a majority-minority district, yet the legislature also formally adopts a partisan-maximization goal that cannot be achieved if that district is preserved, and Black and White voters align strictly with opposite political parties. As Justice Kagan noted in footnote~9 of her dissent (\href{https://supreme.justia.com/cases/federal/us/608/24-109/}{Justia}), \textit{Callais} leaves unresolved whether direct ``smoking-gun'' proof of racial motive allows a Section~2 plaintiff to bypass the new \textit{Gingles} Precondition~1--3 disentanglement requirements, or whether a state's dual partisan motive still defeats liability under \textit{Gingles} while potentially exposing the map to a standalone Fourteenth or Fifteenth Amendment intentional-discrimination claim.
\end{enumerate}

\section{Critique}

\textbf{Originalist, Textualist, and Conservative Perspectives.} Defenders of \textit{Callais}, such as Vanderbilt Law Professor Brian Fitzpatrick, argue that the decision faithfully enforces the Constitution's color-blind equal protection mandate (\href{https://www.scotusblog.com/2026/07/the-case-for-callais/}{SCOTUSblog}). On this view, assigning voters to districts based on skin color is a racial classification subject to strict scrutiny, and achieving ``proportionalish representation'' (racial balancing) has never been a compelling interest in admissions, contracting, or voting. Because the Fifteenth Amendment forbids only intentional racial discrimination, Congress's enforcement power cannot compel states to engage in race-based districting absent a strong inference of purposeful discrimination. Even sympathetic conservative commentators, however, have noted that the majority opinion is awkwardly drafted, declines to acknowledge squarely that it is departing from \textit{Allen v. Milligan}, omits a deep originalist inquiry into the Reconstruction Amendments' original public meaning, and removes one of the last practical checks on aggressive partisan gerrymandering (\href{https://www.scotusblog.com/2026/07/the-case-for-callais/}{Fitzpatrick, SCOTUSblog}).

\textbf{Progressive and Doctrinal Perspectives.} Progressive and election-law scholars---including UC Berkeley Law Dean Erwin Chemerinsky, Harvard Kennedy School analysts, and the Brennan Center for Justice---condemn \textit{Callais} as the judicial nullification of Section~2 (\href{https://www.scotusblog.com/2026/05/understanding-the-recent-voting-rights-act-case/}{Chemerinsky, SCOTUSblog}; \href{https://www.hks.harvard.edu/faculty-research/policy-topics/fairness-justice/what-louisiana-v-callais-means-voting-rights-act}{Harvard Kennedy School}; \href{https://www.brennancenter.org/our-work/analysis-opinion/after-louisiana-v-callais-heres-proof-just-how-bad-voting-rights-america}{Brennan Center}). Chemerinsky situates \textit{Callais} as the culmination of five ideologically divided precedents---\textit{Bolden}, \textit{Shaw v. Reno}, \textit{Shelby County}, \textit{Rucho}, and \textit{Alexander}---that together give zero deference to Congress's 1982 legislative compromise and allow any state to immunize racial vote dilution simply by declaring a partisan gerrymander (\href{https://www.scotusblog.com/2026/05/understanding-the-recent-voting-rights-act-case/}{SCOTUSblog}). Critics also argue that the majority's ``anti-classification'' analogy to university admissions ignores a fundamental distinction: in voting, drawing a remedial district does not exclude any citizen from casting an equally weighted ballot, whereas eliminating Section~2 protection in racially polarized states effectively strips minority communities of equal electoral representation (\href{https://www.scotusblog.com/2026/05/understanding-the-recent-voting-rights-act-case/}{SCOTUSblog}).

\textbf{Underlying Competing Values.} \textit{Callais} crystallizes a profound clash between two visions of constitutional equality and institutional power: an \textit{anti-classification} vision demanding strict color-blindness in state redistricting and cabining federal statutory intrusion into state legislative prerogatives, versus an \textit{anti-subordination} and democratic-representation vision honoring Congress's broad Fifteenth Amendment power to eradicate structural racial vote dilution even when racial and partisan lines overlap.

\section{Key Quotations}

\begin{quote}
``Section 2 of the Voting Rights Act of 1965, 52 U. S. C. \S~10301 \textit{et seq.}, was designed to enforce the Constitution---not collide with it. Unfortunately, lower courts have sometimes applied this Court's \S~2 precedents in a way that forces States to engage in the very race-based discrimination that the Constitution forbids.''
\end{quote}
\noindent --- Justice Samuel A. Alito, Jr., Opinion of the Court, \textit{Louisiana v. Callais}, 608 U.S. \_\_\_ (2026) (slip op., at 1) (\href{https://www.supremecourt.gov/opinions/25pdf/24-109_21o3.pdf}{Slip Opinion}; \href{https://supreme.justia.com/cases/federal/us/608/24-109/}{Justia}; \href{https://www.courtlistener.com/opinion/10852760/louisiana-v-callais-revisions-50426/}{CourtListener}).

\begin{quote}
``In short, \S~2 imposes liability only when the evidence supports a strong inference that the State intentionally drew its districts to afford minority voters less opportunity because of their race. Not only does this interpretation follow from the plain text of \S~2, but it is consistent with the limited authority that the Fifteenth Amendment confers.''
\end{quote}
\noindent --- Justice Samuel A. Alito, Jr., Opinion of the Court, \textit{Louisiana v. Callais}, 608 U.S. \_\_\_ (2026) (slip op., at 26) (\href{https://supreme.justia.com/cases/federal/us/608/24-109/}{Justia}).

\begin{quote}
``This Court should never have interpreted \S~2 of the Voting Rights Act of 1965 to effectively give racial groups `an entitlement to roughly proportional representation.' ... Today's decision should largely put an end to this `disastrous misadventure' in voting-rights jurisprudence.''
\end{quote}
\noindent --- Justice Clarence Thomas, Concurring, \textit{Louisiana v. Callais}, 608 U.S. \_\_\_ (2026) (slip op., at 1) (\href{https://supreme.justia.com/cases/federal/us/608/24-109/}{Justia}).

\begin{quote}
``Under the Court's new view of Section 2, a State can, without legal consequence, systematically dilute minority citizens' voting power. Of course, the majority does not announce today's holding that way. Its opinion is understated, even antiseptic. The majority claims only to be `updat[ing]' our Section 2 law, as though through a few technical tweaks. But in fact, those `updates' eviscerate the law, so that it will not remedy even the classic example of vote dilution given above.''
\end{quote}
\noindent --- Justice Elena Kagan, Dissenting, \textit{Louisiana v. Callais}, 608 U.S. \_\_\_ (2026) (slip op., at 4) (\href{https://supreme.justia.com/cases/federal/us/608/24-109/}{Justia}).

\begin{quote}
``Today's decision renders Section 2 all but a dead letter. In the States where that law continues to matter---the States still marked by residential segregation and racially polarized voting---minority voters can now be cracked out of the electoral process.''
\end{quote}
\noindent --- Justice Elena Kagan, Dissenting, \textit{Louisiana v. Callais}, 608 U.S. \_\_\_ (2026) (slip op., at 47) (\href{https://supreme.justia.com/cases/federal/us/608/24-109/}{Justia}).

\end{document}
